\section{Setting and main contributions}
\label{sec:main:setting}
 Let $\mathcal{D}$ be the set of data $\{\vec{z}_{\nu} \in {\mathbb R}^d, y_{\nu}\}_{\nu \in [n]}$. The input data $\vec{z}_{\nu}$ are taken as a standard i.i.d. Gaussian, while the labels are generated by a teacher, or target, function $y_\nu = f^\star(\vec z_\nu)$. We consider multi-index target function, dependent on a low-dimensional subspace of the input space:
 \looseness=-1 
\begin{align}
   y_\nu = f^\star(\vec z_\nu) = g^\star\left(\frac{W^\star \vec z_\nu}{\sqrt{d}}\right),\, \vec z_\nu \sim {\cal N}(0,{\mathbbm 1}_d)\,,
   \label{eq:main:teacher_def}
\end{align}  We assume for convenience that $W^\star = \{\vec w^\star_r\}_{r \in [k]} \in \mathbb{R}^{k \times d}$ is normalized row-wise on the sphere $\mathcal{S}_{d-1}(\sqrt{d})$, with orthogonal weights i.e $\langle \vec w^\star_l,\vec w^\star_m \rangle = 0$ for $l \neq m \in [r]$. 

The data are handled to a {\it two-layer} network (the student)~$f$ with first layer weights $W = \{\vec w_i \}_{i \in [p]} \in \mathbb{R}^{p \times d}$ ($p$ is the number of neurons in the hidden layer) and second layer weights $\vec a \in \mathbb{R}^p$
with an activation function $\sigma$, that is:
\begin{align}
    f\left( \frac{W\vec z} {\sqrt{d}}\right) = \sum_{j=1}^p a_j \sigma{\left(\frac{\langle \vec w_j, \vec z \rangle}{ \sqrt{d}}\right)}
\end{align}

Our main goal is to analyze the dynamics of gradient descent that minimizes the empirical Mean Squared Error (MSE) loss $\mathcal{L}$ at time $t \in [T]$:
\begin{align}
  {\cal R}_{\rm emprical} &=\sum_{\nu=1}^n\mathcal{L}\left(\frac{W^{(t)}\vec z_\nu} {\sqrt{d}},f^\star(\vec z_\nu)\right) \\
  &= \frac{1}{2} \sum_{\nu=1}^{n}\left[g^\star\left(\frac{W^\star \vec z_\nu}{\sqrt{d}}\right) - f\left( \frac{W^{(t)}\vec z_\nu} {\sqrt{d}}\right)\right]^2\,, \nonumber 
\end{align}
in the high-dimensional limit where $n,d \!\to\! \infty$ with $n/d \!=\! \alpha \!=\! \Theta(1)$. We use a common assumption that is amenable to rigorous theoretical guarantees: we keep the second layer weights $\vec a$ fixed at initialization. 
For convenience, we further impose the constraint of symmetric initialization common in such analyses \citep{dandi2023how,damian2022neural}. Concretely, we assume that the number of neurons $p$ is even and the weights satisfy at initialization:
\begin{equation}\label{eq:init_sym}
    a_i = -a_{p-i+1}, \quad \vec{w}_i^0 = \vec{w}_{p-i+1}^0 \quad \text{for all } i \in [p/2],
\end{equation}
which ensures that the output $f\left( \frac{W^{(t)}\vec z_\nu} {\sqrt{d}}\right)$ equals $0$ at initialization.
For $i, \in [p/2]$, the weights are initialized as $a_i \sim \frac{1}{p}\mathcal{N}(0,1), \vec w_i^{(0)}\!\!\sim\!\! \mathcal{N}(0,{\mathbbm 1}_d)$
Subsequently, with $\vec{a}$ fixed, the first layer weights $W=\{\vec w_i\}_{i \in [p]}$ are learned using gradient descent, producing the following sequence of iterates up to a final time $T$:\looseness=-1
\begin{align} 
    \label{eq:GD_update}
    \vec w_i^{(\!t+1\!)}\!=\!(1\!-\!\eta\lambda)\vec w^{(\!t\!)}_i \!-\! \eta \!\sum_{\nu=1}^n\!\nabla_{\!\!\vec w^{(\!t\!)}_i}
    %_{\vec w^{(t)}_i} 
    \mathcal{L}\!\left(\!\frac{W^{(t)}\vec z_\nu} {\sqrt{d}},f^\star(\vec z_\nu)\!\!\right)
\end{align}
where $\eta\!\in\!\mathbb{R}$ is the learning rate and $\lambda \in \mathbb{R}$ is the explicit regularisation. We may refer to these steps as the \textit{representation learning} steps, in which the first layer weights learn how to adapt to the low dimensional structure identified by the teacher subspace $W^\star$.  
% \\ \lp{To include only if generalization is studied} In this manuscript we do not rigorously specify the implications of the representation learning step's efficiency — or lack thereof — when attempting to learn as well the non-linearity $g_\star$. 

Our \textbf{main contributions} in this paper are the following: 
\vspace{-0.3cm}
\begin{itemize}[noitemsep,leftmargin=1em,wide=0pt]
    \item We characterize the class of multi-index targets that can be learned efficiently by two-layer networks trained with a finite number of iterations of gradient descent in the high dimensional limit $(d \to \infty)$ with large batch sizes ($n = \alpha d, \alpha = O(1)$). We establish a strong separation between what can be learned with one-pass algorithms (that use new fresh batches at every step) and multi-pass gradient approaches that can 
    use the same batch many times (see Figs.     \ref{fig:main:single_index} and \ref{fig:main:multiple_index} for examples).
    \item We show that while both gradient flow \citep{bietti2023learning} and single-pass
    algorithms suffer from the curse of the information exponent \citep{arous2021online}, and are limited to staircase learning \citep{abbe2023sgd}, requiring a diverging number of iterations for non-staircase functions, some of these problems become {\it trivial} when allowing reusing samples multiple times, and features can be learned in just $T=2$ iterations. This disproves, in particular, a recent conjecture by \citep{abbe2023sgd}.
    \item The simplest examples of directions that cannot be learned in a finite number of steps relate to symmetries in the target function. This includes phase retrieval \citep{maillard2020phase} or the specialization transition in committees, as discussed in the Bayes optimal approaches of single-index \citep{barbier2019optimal} and multi-index \citep{aubin2019committee} models.  
    \item The proof of our results is based on the concept of ``hidden progress", and crucially uses the rigorous Dynamical Mean Field Theory (DMFT) \citep{celentano2021high,gerbelot2022rigorous}. This has an interest on its own as it provides a sharp example of how DMFT can help to understand batch reusing to go beyond the current state-of-the-art results.
    %\citep{mignacco2021stochasticity,mignacco2020dynamical}.    
    \item Finally, we use DMFT to provide a closed-form description of the dynamics of gradient descent for two-layer nets. Keeping track of the correlations induced by re-using the same batch leads to a set of integro-differential equations.  We provide rigorous theoretical guarantees in the correlated samples regime without assuming the resampling of a fresh new batch for each iteration of the algorithm. We corroborate the theoretical claims with numerical simulations (See  \href{https://github.com/IdePHICS/benefit-reusing-batch}{\href{https://github.com/IdePHICS/benefit-reusing-batch}{https://github.com/IdePHICS/benefit-reusing-batch}.
}).  
\end{itemize}
\vspace{-0.5cm}



\paragraph{Other Related works --} A major issue in machine learning theory is figuring out how well two-layer neural networks adapt to low-dimensional structures in the data. Different results have tightly characterized the limitations of networks in which the first layer of weights $W$ is kept fixed, i.e. equivalent to kernel approaches \citep{dietrich1999statistical,ghorbani2019limitations,ghorbani2020when,bordelon2020spectrum,loureiro2021learning,cui2021generalization}. This class of learning algorithms, although amenable to theoretical analysis, is unable to learn features in the data. Therefore, one central avenue of research in this context is to understand the efficiency of the \textit{representation learning} (or feature learning) when training with gradient-based algorithms to overcome the limitations of the kernel regime. Sharp separation results between the performance of neural networks at initialization (random features) and trained with only one step of gradient descent (with a large learning rate) have been offered \citep{ba2022high,damian2022neural,dandi2023how}. 

The class of features efficiently learned with multiple steps of one-pass SGD with one sample per batch is characterized by the \textit{information exponent} ($\rm{IE}$) \citep{arous2021online} of the target function. In the context of single-index learning, denoting $\ell$  the $\rm{IE}$ of the target, the algorithm needs $T\!=\! O(d^{\ell -1})$  steps to perform weak recovery of the teacher direction, i.e., obtaining an overlap between learned weights and $\vec w^\star$ better than random guessing \citep{arous2021online}. Recently, these results have been improved up to the Correlational Statistical Query (CSQ) lower bound of $d^{{\max(\frac{\ell}{2}},1)}$, by considering an appropriate smoothing of the loss \citep{damian2023smoothing}. A generalization to large batch one-pass SGD is in \citep{dandi2023how}. 

Similarly, multi-index feature learning presents an unavoidable computational barrier for one-pass algorithms. \citep{abbe2021staircase} first characterizes a hierarchical picture of learning in the Boolean data case: informally, the features efficiently learned at each step of the one-pass algorithm need to be linearly connected with the previously learned features. This concept is formalized by the definition of the staircase property \citep{abbe2021staircase}. This hierarchical picture of learning is extended to large batches in the SGD and non-Boolean data in \citep{abbe2022merged,abbe2023sgd,dandi2023how}. Moreover, \citep{abbe2023sgd} conjecture that re-using the batch can reduce the sample complexity of the target with leap $\ell$ only up to $O(d^{\max(\frac{\ell}{2},1)})$, corresponding to the lower bound for Correlational Statistical Query (CSQ) algorithms. 


We disprove this conjecture and show that the sample complexity for a large class of functions can be reduced to $O(d)$ {\it independently} of the leap exponent $\ell$. More generally, our results show that CSQ lower bounds and the notions of staircase property and information exponent are limited to online-SGD on Gaussian/Boolean data,  and do not describe the class of functions inherently easy or hard to learn by gradient-based methods. We also show that learning non-even single-index functions does not require techniques such as spectral warm-start \citep{chen2020learning}.\looseness=-1

Dynamical Mean Field Theory has a long history in statistical physics. Early theories of dynamics in complex systems were pioneered in soft spin glass models \citep{sompolinsky1981dynamic} and toy models of random feature deep networks \citep{sompolinsky1988chaos}. The DMFT approach used in this paper was first proposed as a way to study ``hard spins" in spin glass models \citep{eissfeller1992new,eissfeller1994mean}, and was later generalized to ``soft spins" \citep{cugliandolo2003dynamics} and more realistic models in condensed matter \citep{georges1996dynamical}. In the context of learning, DMFT was used for optimization problems \citep{mannelli2019who,mannelli2019passed,mannelli2020marvels,mannelli2021just} and for analyzing the behavior and the noise of gradient-based algorithms \citep{mignacco2020dynamical,mignacco2021stochasticity,mignacco2022effective}.  From the mathematics point of view, these DMFT equations were first proven rigorously in the seminal work of \citep{benarous1997symmetric} in the context of spin glasses. Important progress was achieved recently with rigorous proofs of the DMFT equations for multi-index models \citep{celentano2021high,gerbelot2022rigorous} that we use to prove our main results.\looseness=-1

%, and for detailed studies of GD  with a finite learning rate \citep{mignacco2022effective}.

% \et{Would you like to have also something mentioning we do a different DMFT from Pelevan and maybe also Saxe's work?}
% % \citep{cugliandolo1993analytical,benarous1997symmetric,maimbourg2016solution,franz2007dynamics}
% \\ \et{not} 





%%% PREVIOUS %%%
% \paragraph{Sufficient statistics}
% Our study, like many other efforts \citep{arous2022high,saad1995line}, is based on the concentration of the neurons' overlaps with the target subspace and their norms, rather than the concentration of the complete gradient vectors. This approach only requires the knowledge of the following \textit{sufficient statistics} determining the dynamics; let the pre-activations be defined as:
% \begin{align}
%     \vec \lambda = \vec W z \qquad  \text{and} \qquad \lambda^\star = W^\star \vec z
% \end{align}
% Thanks to the Gaussian nature of the data, the pre-activations are jointly Gaussian vectors $(\vec{\lambda}^{\nu}, {\vec{\lambda}^{\star}}^{\nu})\sim\mathcal{N}(\vec{0}_{p+k}, \Omega)$ with covariance: 
% \begin{equation}
% \Omega \coloneqq
% \begin{pmatrix}
% Q & M\\
% {M^{\top}} & P
% \end{pmatrix}=
% \begin{pmatrix}
% W{W}^{\top} & W{W^{\star}}^{\top}\\
% W^{\star}W^{\top} & W^{\star} W^{\star T}
% \end{pmatrix}
% \in\mathbb{R}^{(p+k)\times (p+k)}
% \end{equation}
% The theoretical analysis is based on the description of the evolution of the covariance matrix $\Omega$ as a function of time.
